% !PRINTED 117
\weberpage{117}
%<<CONT
cyklisch ($S = 1, \gamma, \gamma^2, \gamma^3, \gamma^4, \gamma^5$). Es bleibt also nur die Annahme, dass $\gamma$ vom $2^{\mathrm{ten}}$ Grade, also
\weq{2}{\gamma^2 = 1}
ist. Da wir aber $\gamma$ durch $\gamma\alpha$ oder $\gamma\alpha^2$ ersetzen können, so müssen auch diese vom $2^{\mathrm{ten}}$ Grade sein und es folgt:
\weq{3}{\gamma\alpha = \alpha^2\gamma, \quad \gamma\alpha^2 = \alpha\gamma,}
mit deren Hülfe man jedes Compositum aus beliebig vielen Potenzen von $\alpha$ und $\gamma$ immer auf eines und nur eines der Elemente $S$ zurückführen kann. Und wenn man jetzt
\[
1, \alpha, \alpha^2, \gamma, \gamma\alpha, \gamma\alpha^2
\]
mit
\[
1, \alpha, \beta, \gamma, \delta, \varepsilon
\]
bezeichnet, so erhält man folgende Tabelle:
\[
\begin{array}{c!{\vrule width 0.8pt}c|c|c|c|c|c}
 & 1 & \alpha & \beta & \gamma & \delta & \varepsilon \\
\noalign{\hrule height 0.8pt}
1 & 1 & \alpha & \beta & \gamma & \delta & \varepsilon \\ \hline
\alpha & \alpha & \beta & 1 & \varepsilon & \gamma & \delta \\ \hline
\beta & \beta & 1 & \alpha & \delta & \varepsilon & \gamma \\ \hline
\gamma & \gamma & \delta & \varepsilon & 1 & \alpha & \beta \\ \hline
\delta & \delta & \varepsilon & \gamma & \beta & 1 & \alpha \\ \hline
\varepsilon & \varepsilon & \gamma & \delta & \alpha & \beta & 1
\end{array}
\]

Es kommen, wie es sein muss, in jeder Zeile und in jeder Colonne alle Elemente vor.

\begin{center}
\S.\,28.

\gesperrt{Beziehung der allgemeinen Gruppen zu den Permutationsgruppen.}
\end{center}

Die Gruppen, die wir bisher am meisten angewandt haben, sind die Permutationsgruppen; diese Art von Gruppen gewinnen eine erhöhte Bedeutung auch für die allgemeine Gruppentheorie durch die folgenden Betrachtungen.

Es sei
\weq{1}{P = a_0, a_1, a_2, \ldots, a_{n-1}}
eine beliebige Gruppe vom Grade $n$. Greifen wir aus $P$ irgend
%CONT>>
